\section{A complete segment trace estimate}
\label{sec:v87-trace}
The additional optical moment can also be used without selecting a
large-sensitivity sector. We apply a one-dimensional trace argument to
an exact physical angular density, and sum the resulting energies through
the Borel segment bijection. The power strictly above one is essential.

Fix $1<p<17/16$ and put $e_p=1-1/p$. Write
\[
 u_p(t)\,dt=F_\#(I_m^p\bar\nu),\qquad
 M_p(J)=\int_Ju_p(v)\,dv.
\]
The endpoint gradient has a uniform upper bound $z\le z_{\max}$.
Consequently $I_m^p\le z_{\max}^{\eta}I_m^pz^{-\eta}$ on the
regular source, for any fixed $0<\eta<1-p/2$. The actual mixed-moment
Theorem~\ref{thm:v86-moments} therefore implies
\begin{equation}\label{eq:v87-moment}
 m^2M_p(J)\le C_p(|J|+B_*^{-1})
  +C_{B_*,p}(1+B_*|J|)m^2\rho_{B_*}^{m(1-p/2)/2}.
\end{equation}
This is a finite-count bound for every interval, not just an assertion
about one previously fixed window. Its spectral inputs are exactly those
of the retained mixed-moment theorem.

\begin{lemma}[A positive-density trace inequality]
\label{lem:v87-positive-trace}
Let $A$ be positive and locally absolutely continuous on an interval,
and let $I\ge1$ be measurable. Suppose that $|A'|\le DAI$.
Let $J_0$ be an interior interval of length $d>0$, and let $K$ be a
set of target points such that the intervals joining $J_0$ to those
points lie in an interval $J_1$ of length at most $8d$. Then
\begin{equation}\label{eq:v87-positive-trace}
 \sup_K A\le C_p\{d^{-1}+D^pd^{p-1}\}\int_{J_1}AI^p.
\end{equation}
The inequality also holds for one-sided limits of target points when
the right side is finite.
\end{lemma}
\begin{proof}
Set $f=A^{1/p}$. On every compact interior interval,
$|f'|\le (D/p)fI$. There is a point $v_0\in J_0$ at which
$f(v_0)\le(d^{-1}\int_{J_0}A)^{1/p}$, or a sequence of points with
this bound up to an arbitrarily small error. The fundamental theorem
of calculus and H\"older's inequality give, for $v\in K$,
\[
 f(v)\le \left(d^{-1}\int_{J_0}A\right)^{1/p}
       +\frac Dp(8d)^{1-1/p}\left(\int_{J_1}AI^p\right)^{1/p}.
\]
Raise to the power $p$, use $(x+y)^p\le2^{p-1}(x^p+y^p)$ and
$I\ge1$, and exhaust the joining intervals by compact subintervals.
This proves the assertion. No bound for $A^{-1}$ is used.
\end{proof}

For $0<a,\delta\le1$, define the following complete initial family
on the roof $t$:
\begin{equation}\label{eq:v87-initial-family}
 \mathcal E_t(a,\delta)=
 \{(j,\theta):a_j(\theta)<t<b_j(\theta),\
                  z_j(t,\theta)\ge a,\ b_j(\theta)-a_j(\theta)\ge\delta\}.
\end{equation}
All endpoint types are included. Here and below $a$ without a subscript
is a gradient threshold, and $a_j$ is an endpoint roof. Set
\begin{equation}\label{eq:v87-trace-width}
 d=d(a,\delta)=\min\{\delta/16,\ a^2/(32\Lambda_\circ),\ 1/16\},
 \qquad J_t=(t-4d,t+4d).
\end{equation}
For the physical constants in this article, $d\le a$. The source
selection is a restriction of the original exact event, not a new
normalization or a replacement of the full source.

\begin{theorem}[Summed physical trace with multiplicity one]
\label{thm:v87-summed-trace}
For almost every $t$ at each fixed count and exact parameter set,
\begin{equation}\label{eq:v87-summed-trace}
 S_t(a,\delta):=
 \sum_j\int_{\mathcal E_t(a,\delta)\cap(\{j\}\times E_j)}
    \sup_{t\le v<\min(t+d,b_j)}A_j(v,\theta)\,d\theta
                  \le C_p d^{-1}M_p(J_t).
\end{equation}
The same bound holds for the original unweighted source density
restricted to $z\ge a$ and total regular residence at least $\delta$.
In particular, for every fixed positive $a,\delta$,
\begin{equation}\label{eq:v87-collared-source-height}
 \limsup_m\sup_{R,n,k}m^2\|p^{a,\delta}\|_\infty\le C_p,
\end{equation}
where the leading constant is independent of those fixed cutoffs.
Equation~\eqref{eq:v87-collared-source-height} does not remove them.
\end{theorem}
\begin{proof}
At least one side of a selected segment has length at least
$\delta/2\ge8d$ from $t$. Use the anchor interval
$(t-3d,t-2d)$ on the backward side when that side is long, and
$(t+2d,t+3d)$ otherwise. This is a Borel choice. Every path from
an anchor to a target in \eqref{eq:v87-summed-trace} remains in the
same physical segment and in $J_t$. Radial growth gives there
\[
 z_j(v)^2\ge a^2-8\Lambda_\circ d\ge3a^2/4.
\]
The established inequality $|\partial_vA_j|\le C_0A_jI_m/z_j$
thus permits Lemma~\ref{lem:v87-positive-trace} with $D=2C_0/a$.
Since $d\le a$, its coefficient is at most $C_pd^{-1}$.

Integrate over the selected angular parameters. At every interior
roof the Borel bijection counts each physical point once, including
its exact labels. Therefore the sum of the energies on all joining
intervals is bounded by $M_p(J_t)$, with no word-count factor.
Suprema over variable target intervals are measurable: continuity on
each regular segment permits a countable dense set of roof offsets.
The source moment is finite, so truncation and positive convergence
justify this integration also for countably many segment charts.
The common physical kernels fix the almost-everywhere convention.

The density $p^{a,\delta}(t)$ is bounded by $S_t$. Insert
\eqref{eq:v87-moment}, fix $a,\delta,B_*$, take the collision
limsup, and only then enlarge $B_*$. The length of $J_t$ is $8d$,
which cancels $d^{-1}$. This proves
\eqref{eq:v87-collared-source-height}. The finite-count estimate still
contains the explicit $d^{-1}B_*^{-1}$ and spectral remainder;
it does not imply the height bound at zero cutoffs.
\end{proof}
